\documentclass[11pt]{article}
\usepackage[a4paper,margin=25mm]{geometry}
\usepackage{amsmath,amssymb,amsthm,booktabs,longtable,xurl,hyperref}
\renewcommand{\labelitemi}{\(\bullet\)}
\hypersetup{hidelinks}
\newtheorem{theorem}{Theorem}
\newtheorem{lemma}[theorem]{Lemma}
\newtheorem{proposition}[theorem]{Proposition}
\theoremstyle{remark}\newtheorem{remark}[theorem]{Remark}
\newcommand{\Q}{\mathbb Q}
\newcommand{\Tr}{\operatorname{Tr}}
\newcommand{\Nm}{\operatorname{N}}
\title{Removing private scalar normalizers from a common point chart}
\author{Z diophantine in Q project}
\date{23 November 2026}
\begin{document}
\maketitle
\begin{abstract}
We give three exact changes to the complete quadratic systems with a
common contact normalization. The changes replace private multiples of
the same scalar point by common coordinates and remove reciprocal
normalizers that the new coordinates make unnecessary. The resulting
systems have 1143 variables and 987 quadratic equations in the general
family, and 1323 variables and 1143 quadratic equations in the specified
coefficient family. These counts precede the separate nine-variable
Pell reduction. Every coefficient table is materialized, and an
independent checker reconstructs every row and verifies the recovery
identities. This paper proves equivalence to the preceding systems;
it does not prove their unresolved source universality hypothesis.
\end{abstract}

\section{Statement and scope}

The preceding common-normalization construction has twelve rational
integrality blocks, one for each input of the retained emended Jones
system. Its two complete tables have respectively 1185 variables and
1095 equations, or 1355 variables and 1241 equations
\cite{global}. The latter table uses the specified effective integrality
interface; the former belongs to the family with an ineffective
sufficient threshold. In both tables, each residual is an ordinary
quadratic polynomial with integer coefficients in the fixed parameters.

The source hypothesis denoted by \(U_{\rm em}\) is the universality of
the precise emended source recorded in the repository. It has not been
independently established. In particular, its printed scalar parameter
still contains \(n^2\); this paper makes no correction to that formula
and does not silently substitute a different source. The arithmetic
interfaces stated in \cite{global,scalar} are also retained. No new
\(abc\) assumption is introduced. These qualifications apply equally
before and after the changes proved here.

\begin{theorem}\label{thm:main}
There are exact changes of auxiliary coordinates between the complete
common-normalization systems and complete quadratic systems with the
following counts:
\begin{center}
\begin{tabular}{lrrrr}
\toprule
Interface & Old variables & New variables & Old equations & New equations\\
\midrule
General, 163 coordinates &1185&1143&1095&987\\
Specified, 179 coordinates &1355&1323&1241&1143\\
\bottomrule
\end{tabular}
\end{center}
For every admissible fixed parameter choice, the changes preserve the
rational solution projection to all source inputs. Each new system therefore has a sum of squares
of quadratic residuals of degree exactly four, in the stated number of
rational variables. All arithmetic and universality hypotheses of the
predecessor are preserved without strengthening.
\end{theorem}

The independent Pell-core reduction removes another nine coordinates
and three equations \cite{pell}. It concerns different coordinates and
rows. Applying it after this theorem gives, respectively, 1134 variables
with 984 equations, and 1314 variables with 1140 equations. Those
combined counts are not already included in the two tables of this
paper. Neither statement asserts a global minimum, nor undecidability
under hypotheses that have not been proved.

\section{The common and private quantities}

Work over the field \(K=\Q(\sqrt2)\) and put
\[
                       D=11-6\sqrt2.
\]
A field coordinate is represented by its two rational components. If
\(a=a_0+a_1\sqrt2\) and \(b=b_0+b_1\sqrt2\), multiplication is
\[
 ab=(a_0b_0+2a_1b_1)+(a_0b_1+a_1b_0)\sqrt2.
\]
Thus a quadratic equation over \(K\) gives two ordinary rational
quadratic equations.

Let \(i=1,\ldots,k\) index the integrality blocks; the complete
application uses \(k=12\). Each block has its own positive rational
normalizer \(R_i\), field coordinate \(\widehat J_i\), and rational
coordinate \(\widehat\tau_i\). The fields
\[
        f,\quad S,\quad Z,\quad U,\quad V,\quad W,
        \quad F_U,\quad F_V,\quad A
\]
are already common in the preceding construction. The square \(f^2\)
is stored there. Write
\[
                 L_i=\widehat J_i-f^2,
                 \qquad M=-2S-W.
\]
The scalar equations are
\begin{equation}\label{eq:scalar}
       R_iL_i=Z^2,\qquad
       \widehat\tau_i\Tr(L_i)=\Nm(\widehat J_i).
\end{equation}
The field \(L_i\) and \(M\) are expressions, not extra coordinates.

The old scalar point chart uses private fields \(P_i^*\) and
\(r_i^*\), with a common field \(r_2\), and satisfies
\begin{equation}\label{eq:old}
\begin{aligned}
 UV&=L_iP_i^*,& P_i^*W&=R_i r_i^*,\\
 r_i^*(U+V-W-4S)+DR_iL_i&=0,&
 L_iP_i^*+W(U+V)&=M^2,\\
 r_2&=(r_i^*)^2,& F_Ur_2&=U^2,\\
 F_Vr_2&=V^2,& A(U+M)&=2SZ.
\end{aligned}
\end{equation}
These are physical coordinates, before their monomial rescalings in
the stored rational tables. The point construction in \cite{scalar}
identifies \(P_i^*=R_i x(Q_2)\); we shall not need to assume this
interpretation in order to prove the algebraic sharing.

The stored general table uses \(\widetilde r_i=R_i r_i^*\), so its
copies of the third and fifth equations in \eqref{eq:old} are
\begin{equation}\label{eq:stored-r}
\begin{aligned}
 \widetilde r_i(U+V-W-4S)+D R_{2,i}L_i&=0,\\
 R_{2,i}r_2-\widetilde r_i^2&=0,
 \qquad R_{2,i}=R_i^2.
\end{aligned}
\end{equation}
The specified table already stores \(r_i^*=r\) as a common field.
Its stored square is \(r_2/c\), where the nonzero global normalizer
\(c\) is common. Consequently its square equation is
\(c(r_2/c)=r^2\), still quadratic. No change to this convention
is needed.

\section{Nonvanishing on the full rational solution set}

\begin{lemma}\label{lem:nonzero}
In both the old and new systems, \(R_i\ne0\), \(f\ne0\), and
\(Z\ne0\) on every rational solution. These facts hold before
applying an integrality or universality conclusion.
\end{lemma}
\begin{proof}
The retained rational normalizer has, for a rational input \(X\),
the form
\[
 \beta=51X^2+2504,\qquad
 r_X=\beta(\beta-1)(\beta-2401)(32\beta+343)>0.
\]
In the general family,
\[
 R_i=\frac{149C^2r_X}{1+149^2r_X},\qquad C=149^N\ne0;
\]
in the specified family,
\[
 R_i=\frac{A_{\rm a}r_X}{1+25r_X},\qquad A_{\rm a}>0.
\]
The reciprocal input block uses its retained guarded rational input;
the same argument applies. The common reference chart explicitly
guards the finite nonzero reciprocal \(f\), as in
\cite{scalar,global}.

If \(Z=0\), the first equation of \eqref{eq:scalar} implies
\(L_i=0\), and hence \(\widehat J_i=f^2\). The second then gives
\(\Nm(f^2)=0\), impossible in the field \(K\) when \(f\ne0\).
In the new system, the equation \(K_0=Z^2\) and
\(R_iL_i=K_0\) give the identical argument. Here \(K_0\) is the
field square coordinate introduced below, not the ground field.
\end{proof}

\section{One common scalar abscissa}

Introduce common fields
\[
                         K_0=Z^2,\qquad p,
\]
and, in the general table, a common field \(r\). Replace the first
five equations in \eqref{eq:old}, together with the first equation
in \eqref{eq:scalar}, by
\begin{equation}\label{eq:new}
\begin{aligned}
 K_0&=Z^2,& R_i(\widehat J_i-f^2)&=K_0,\\
 UV&=K_0p,&pW&=r,\\
 r(U+V-W-4S)+DK_0&=0,&UV+W(U+V)&=M^2,\\
 r_2&=r^2.
\end{aligned}
\end{equation}
Only the equation containing \(R_i\) is repeated for each block.
Retain the trace equation, the last three equations of
\eqref{eq:old}, and every other unchanged source equation.
Each equation in \eqref{eq:new} is quadratic in the stored rational
coordinates. In the specified table, its last equation is written
\(c(r_2/c)=r^2\).

\begin{proposition}\label{prop:rechart}
For every \(k\ge1\), equations \eqref{eq:scalar} and
\eqref{eq:old} are equivalent to the stated new equations, with
recovery maps
\[
       P_i^*=R_ip,\qquad r_i^*=r,\qquad
       \widetilde r_i=R_ir.
\]
The equivalence holds on the full solution sets satisfying the
retained scalar and reference equations.
\end{proposition}
\begin{proof}
Start with an old solution. By Lemma~\ref{lem:nonzero},
\(Z^2\ne0\) and \(R_i\ne0\). From \(R_iL_i=Z^2\) and
\(UV=L_iP_i^*\), we obtain
\[
                 \frac{P_i^*}{R_i}=\frac{UV}{Z^2}.
\]
The right side is independent of \(i\). Define it to be \(p\).
The old equation \(P_i^*W=R_i r_i^*\) then gives
\(r_i^*=pW\); all these fields have the same value \(r\).
Substituting in the remaining equations gives \eqref{eq:new}.
The changes to the aggregate contributions are
\[
          P_i^*\longmapsto R_ip\quad\hbox{in the general family},
          \qquad \frac{P_i^*}{R_i}\longmapsto p
             \quad\hbox{in the specified family}.
\]
The former is quadratic and the latter linear.

Conversely, take a new solution. Define \(P_i^*=R_ip\) and
\(r_i^*=r\), and use \(\widetilde r_i=R_ir\) in the general
stored table. Then \(L_iP_i^*=K_0p=UV\) and
\(P_i^*W=R_ir\). The remaining old equations follow directly
from \eqref{eq:new} and the retained equality \(R_{2,i}=R_i^2\).
The aggregate substitutions are exact identities.

Only \(R_i\) and \(Z^2\), whose nonvanishing is already proved,
were divided out. No secant denominator, alternate-point ordinate,
or native-pole auxiliary was inverted. Finite tangent branches and
allowed alternate two-torsion branches therefore remain present.
For completeness, the new sum equation forces \(r\ne0\), since
\(D K_0\ne0\). Thus \(pW=r\) forces \(p,W\ne0\), and
\(UV=K_0p\) forces \(U,V\ne0\). No additional zero component of
the point chart has been introduced.
\end{proof}

The proof admits a useful general form: whenever a collection of
quadratic blocks contains \(R_iL_i=K_0\), \(UV=L_iP_i\), and
the common \(K_0\) is nonzero, the quotients \(P_i/R_i\) are
already forced to coincide. Introducing that common quotient can
remove the private multiplication registers while retaining degree
two. This is stronger than choosing simultaneous witnesses that
happen to agree: the agreement holds on every solution of the old
system.

\section{Removing twelve inverse normalizers}

The specified predecessor has an extra coordinate
\(R_i^{-1}\) and its guard \(R_i R_i^{-1}=1\). After
Proposition~\ref{prop:rechart}, that inverse disappears from both
aggregate components. Its only remaining substantive occurrence is
in a homogeneous native-pole input equation.

We describe the complete change, so that zero auxiliary branches
are explicitly covered. Use physical native-pole coordinates
\[
           E,U_s,V_s,P_s,Q_s,\sigma,B,P_2,Q_2.
\]
The pole equations can be written
\begin{equation}\label{eq:pole}
\begin{aligned}
 \widehat\tau_i\sigma&=R_{2,i}^{\,2},&
 B\sigma&=R_{2,i}+9792R_i,\\
 \sigma P_2&=P_s^2,&\sigma Q_2&=Q_s^2,\\
 3E^2-34E\sigma-289\sigma^2+17BP_2&=0,\\
 E(17\sigma-B)-17U_sV_s&=0,\\
 U_s^2+34U_sV_s-867V_s^2-17BQ_2&=0.
\end{aligned}
\end{equation}
All rows are quadratic, using the already retained coordinate
\(R_{2,i}=R_i^2\).

The specified predecessor stores each of these nine coordinates
divided by \(R_i\). Denote those old coordinates by primes. Its
first two equations are consequently
\[
 \widehat\tau_i\sigma'=R_iR_{2,i},\qquad
 B'\sigma'=1+9792R_i^{-1}.
\]
The other five equations have exactly the same quadratic shape in
the primed coordinates, because every term has degree two.

\begin{lemma}\label{lem:pole}
Replacing the nine primed coordinates by their unprimed physical
values, using \eqref{eq:pole}, and deleting \(R_i^{-1}\) and its
unit equation preserves the complete rational solution projection.
\end{lemma}
\begin{proof}
The forward map multiplies each primed coordinate by \(R_i\).
Multiply the old first equation by \(R_i\), and the old second
by \(R_i^2\); use \(R_{2,i}=R_i^2\). This gives the first two
equations of \eqref{eq:pole}. Multiplication by \(R_i^2\) gives
the remaining five equations. Conversely divide the nine new
coordinates by the nonzero rational \(R_i\), and restore the
unique coordinate \(R_i^{-1}\). The inverse calculations give
the old system. No native-pole coordinate is assumed nonzero.
In particular a zero value of any of \(E,U_s,V_s,P_s,Q_s,P_2,Q_2\)
is transported without exclusion.
\end{proof}

The general table already has the unprimed form
\eqref{eq:pole}; it needs no pole reweighting. This distinction is
why the third reduction saves twelve variables only in the
specified table.

\section{Coordinate counts and exact verification}

For arbitrary \(k\), the general table loses \(2k\) rational
coordinates for the private \(\widetilde r_i\) and \(2k\) for
the private \(P_i^*\), and adds six for \(r,p,K_0\). Its saving
is \(4k-6\). The specified table already shares \(r\); it loses
\(2k\) coordinates for \(P_i^*\), adds four for \(p,K_0\), and
loses the \(k\) inverse-normalizer coordinates of
Lemma~\ref{lem:pole}. Its saving is \(3k-4\).

For \(k=12\), these are respectively 42 and 32. In the general
case five old field equations per block are replaced by six common
field equations, a saving of \(10k-12=108\) rational rows. In the
specified case four old field equations per block and one common
field square equation are replaced by six common field equations;
the twelve inverse unit equations are also removed. The saving is
\(8k+2-12+k=98\) rows. The private scalar field equations are
modified but their number is unchanged. This proves the counts in
Theorem~\ref{thm:main}.

The fixed source residuals outside the integrality blocks are
untouched. A nonzero square monomial remains in a residual, so
the sum of squared residuals has degree exactly four: its leading
homogeneous part is a sum of squares over the real field and cannot
cancel identically. Over \(\Q\), a sum of rational squares is zero
if and only if every residual is zero.

The reproducible implementation consists of:
\begin{itemize}
\item
\nolinkurl{undecidability/code/build_scalar_rechart.py}, which reads
the frozen complete predecessor tables and writes the new tables;
\item
\nolinkurl{undecidability/verify_scalar_rechart.py}, which independently
reconstructs every retained and changed row in \(\Q[\sqrt2]\),
checks the polynomial recovery identities, audits every occurrence
of each deleted or reweighted coordinate, and verifies the census,
degree, integer coefficient domain, and source-status fields;
\item the complete coefficient tables and hash certificates in
\nolinkurl{undecidability/coefficients/jones1980_scalar_rechart/}.
\end{itemize}
The exact checker additionally rebuilds each table and compares its
canonical serialized bytes. Its input and output hashes prevent a
successful check from being confused with a check of a different
predecessor. These finite checks support the algebraic compilation.
They do not certify the retained arithmetic interfaces or
\(U_{\rm em}\), which remain explicitly identified hypotheses.

\begin{thebibliography}{9}
\bibitem{global}
Z diophantine in Q project,
\emph{One contact gauge in a combined quadratic system},
23 November 2026,
\nolinkurl{undecidability/papers/global_gauge_quartic_checkpoint.tex}.
\bibitem{scalar}
Z diophantine in Q project,
\emph{A scalar-point sharing construction in 163 and 179 variables},
20 November 2026,
\nolinkurl{papers/quartic/scalar_point163_179_20261120_v82.tex}.
\bibitem{pell}
Z diophantine in Q project,
\emph{Pell index optimization}, 23 November 2026,
\nolinkurl{undecidability/papers/pell_index_optimization.tex}.
\end{thebibliography}
\end{document}
